 \documentclass[11pt]{article}

 \input{0-parameters.tex}

\begin{document}


\section{Notes from Fran\c{c}ois, 16 March}

We take as accounting truth that
\[
    \hat{A} = \sum_{i=1}^N \lambda_i \hat{A}_i,
\]
and we define
\[
\hat{\tilde{A}}_i \equiv \left(\zeta\hat{X}_i+(\zeta - 1) \hat{p}_i\right).
\]

We also know that, by definition of $\zeta_i$,
\[
\hat{A}_i \equiv \left(\zeta_i \hat{X}_i+(\zeta_i - 1) \hat{p}_i\right).
\]

In other words, we can write
\begin{equation}
\hat{A}   = \sum_{i=1}^N \lambda_i \hat{\tilde{A}}_i + \sum_{i=1}^N \lambda_i \mathcal{E}_i,
\label{eq:ahat_plus_error}
\end{equation}
where
\begin{equation}
\mathcal{E}_i \equiv  (\zeta_i - \zeta) \left( \hat{X}_i + \hat{p}_i\right).
\label{eq:eps}
\end{equation}
[I worry I made an algebra mistake, to check. ]\ER{None detected by redoing the maths.}


Now, Emilien, you showed below that 
\begin{equation}
\sum_{i=1}^N \lambda_i \hat{\tilde{A}}_i = \sum_{f=1}^F \frac{w_fL_{f}}{pY}\hat{w}_f.
\label{eq:emilien_result}
\end{equation}
One needs to check, but I think in general it must be true that
\begin{equation}
\hat{A} = \sum_{f=1}^F \frac{w_fL_{f}}{pY}\hat{w}_f,
\label{eq:francois_guess}
\end{equation}
since this is just the dual definition of aggregate TFP. Inserting my to-be-proven guess \eqref{eq:francois_guess} and your result \ref{eq:emilien_result} into what we agreed is an accounting definition \eqref{eq:ahat_plus_error}, we get
\[
\sum_{i=1}^N \lambda_i \mathcal{E}_i = 0.
\]
It would be nice to try and prove this separately, starting from \eqref{eq:eps}.

\ER{Proof}\\
\er{Following \cite{hsieh_productivity_1999} and \cite{antras_factor_2003}, we start from the accounting identity that total value added, GDP, $Y$ is equal to the sum of factor payments
\begin{equation}
    Y=\sum_{f=1}^F w_f L_f.
\end{equation}
Differentiating, we get
\begin{equation}
    \hat{Y}=\sum_{f=1}^F \frac{w_fL_{f}}{pY}(\hat{L}_f+\hat{w}_f),
\end{equation}
which we rearrange as
\begin{equation}
    \underbrace{\hat{Y}-\sum_{f=1}^F \frac{w_fL_{f}}{pY}\hat{L}_f}_{\text{Agg. TFP primal}}= \underbrace{\sum_{f=1}^F \frac{w_fL_{f}}{pY}\hat{w}_f}_{\text{Agg. TFP dual}},
\end{equation}
where the left-hand side is the classic definition of the aggregate Solow residual in its primal form, thus equal to the weighted sum of the growth rate of real factor prices, so-called dual.}
\er{\textbf{qed}}

If this works, this would suggest we should really do the backtesting ``jointly'' (industry-level + Domar aggregation)? And not worry about upstream and downstream industries, just sum up the productivities?
\ER{But what about the CCC case where we only have some sectors? Or the ``only energy sector case'': we cannot consider the battery industry separately, it makes sense to only consider the ``downstream'' energy sector there in my opinion.}



\section{Following our discussion on double counting and ``upstream'' vs ``downstream''}

\subsection{Set up}


\paragraph{Case} We are considering in the CCC application (and in later applications to some key sectors) only the key sectors for the net zero. 
But there is a risk that these sectors see decrease in price, not because they're more productive, but because their inputs are being produced more productively.
In a simplified vision:  An upstream industry $U$ and a downstream sector $M$ consuming $U$ as an input. If $M$ appears more productive (i.e. offers a lower price) only because it purchases cheaper inputs from $U$, do we need to model $U$ explicitly to capture the net productivity gains? Or is it enough to consider $M$ only?

I show in the dedicated appendix/document that the correct approach is then to only consider the $M$ sector with our framework. This lead to caveats.


\paragraph{Caveats} 
\begin{enumerate}
    \item We show that we have unbiased estimate of aggregate productivity when we are using all available sectors ($N$). Does this mean we introduce a bias when not considering all sectors? Should we calibrate $\zeta$ to reduce the error between the two sides of the equation $ \hat{A}=\sum_{i=1}^N\lambda_i (\zeta\hat{X}_i+(\zeta-1)\hat{p}_i)$?
    \item There is a risk of double counting if we consider both industries U and M within our framework.
    \item The $\zeta$ is already supposed to take care of the discrepancy between productivity growth and price growth. 
\end{enumerate}

\ER{I try and address the three points successively in this document.}

\subsection{Solution \#1}
\subsubsection{Production network approach}

With the same terminology as in \cite{mcnerneyHowProductionNetworks2022}, where $\mathbf{A}$ includes both intermediate inputs and factors, same for $\mathbf{p}$.

Dual productivity can be given in matrix form,
$   \mathbf{\hat{p}}=-\mathbf{\hat{A}} + \mathbf{A}'\mathbf{\hat{p}} \Rightarrow
 (\mathbf{I}-\mathbf{A}')\mathbf{\hat{p}}=-\mathbf{\hat{A}} $

We can rewrite it as follows, with the assumption that $\zeta_i$ is calibrated exactly on each industry so that $\mathbf{\hat{\tilde{A}}}=\mathbf{\hat{A}}$,
\begin{equation*}
    \mathbf{\hat{\tilde{A}}}=\bm{\zeta} \mathbf{\hat{X}}+(\bm{\zeta}-\mathbf{I})\mathbf{\hat{p}} 
\end{equation*}

Hence in matrix form, we can rearrange,

\begin{align*}
    & \mathbf{\hat{A}}=- (\mathbf{I}-\mathbf{A}')\mathbf{\hat{p}} = \bm{\zeta} \mathbf{\hat{X}}+(\bm{\zeta}-\mathbf{I})\mathbf{\hat{p}} \\
    \Rightarrow& 
    - \left[\mathbf{I}-\mathbf{A}'+(\bm{\zeta}-\mathbf{I})\right]\mathbf{\hat{p}} = \bm{\zeta} \mathbf{\hat{X}} \\
     \Rightarrow &
    - (\bm{\zeta}-\mathbf{A}')\mathbf{\hat{p}} = \bm{\zeta} \mathbf{\hat{X}}.
\end{align*}


One could go further and express $\bm{\zeta}$ as the sum of a common value $\bar{\zeta}$ and an ``error'' term $\Delta \zeta_i$, such as 
\begin{equation*}
    \bm{\zeta}=\bar{\zeta}\mathbf{I}+\bm{\Delta \zeta_i}.
\end{equation*}
The assumption in the paper is that $\bm{\Delta \zeta_i}=0$ on average across sectors and after 20 years.

\vspace{3mm}

\textcolor{orange}{Where do we go from there?}


\subsubsection{Combining the zeta assumption and the Domar aggregation}

We have N sectors in the economy.
Let us assume that $\forall i,~\hat{X}_i=0$ and that the factor prices are constant $\forall f, ~ \hat{w}_f =0$.\textit{ We relax these assumptions in the next section.}

Hence, industry-level productivity prediction is $\hat{\tilde{A}}_i \equiv (\zeta - 1) \hat{p}_i$.
Aggregate productivity is $\hat{A} \equiv\sum_{i=1}^N \lambda_i \hat{A}_i$. 

Let us concatenate the two equations, assuming that $\hat{A}_i= \hat{\tilde{A}}_i$, which mean that $\forall i,~\zeta_i=\zeta$. Note that we prove numerically only that $\hat{A}_i\approx \hat{\tilde{A}}_i$. 

Let us remember that the dual definition of productivity gives $\hat{A}_i \equiv - \left( \hat{p}_i - \sum_{j=1}^N  a_{ji} \hat{p}_j - \sum_{f=1}^F \tilde{l}_{fi} \hat{w}_f \right)$. Which in this case, with $\hat{w}_f =0$, gives $   \hat{p}_i=-\hat{A}_i + \sum_{j=1}^N  a_{ji} \hat{p}_j $.

Implicitly, what we are doing is equating $- \hat{p}_i + \sum_{j=1}^N  a_{ji} \hat{p}_j =(\zeta - 1) \hat{p}_i$. 

\begin{align*}
    \hat{A}&=\sum_{i=1}^N \lambda_i \hat{A}_i\\
    & = \sum_{i=1}^N \lambda_i \hat{\tilde{A}}_i\\
    & = \sum_{i=1}^N \lambda_i (\zeta - 1) \hat{p}_i &\text{(Using that }\forall i, ~\hat{A}_i= \hat{\tilde{A}}_i) \\
    & = \sum_{i=1}^N \lambda_i (\zeta - 1) \left(-\hat{A}_i + \sum_{j=1}^N  a_{ji} \hat{p}_j\right) &\text{(Using the dual def. of }\hat{A}_i) \\
    & = -(\zeta - 1)\sum_{i=1}^N \lambda_i  \hat{A}_i +\sum_{i=1}^N(\zeta - 1) \lambda_i \left(\sum_{j=1}^N  a_{ji} \hat{p}_j\right)\\
    & = -(\zeta - 1)\sum_{i=1}^N \lambda_i  \hat{A}_i +\sum_{i=1}^N(\zeta - 1)  \left(\sum_{j=1}^N  \lambda_i a_{ji} \hat{p}_j\right)\\
        & = -(\zeta - 1)\sum_{i=1}^N \lambda_i  \hat{A}_i +\sum_{i=1}^N(\zeta - 1)  \left(\sum_{j=1}^N  \frac{p_iX_i}{pY} \frac{p_jX_{ij}}{p_iX_i} \hat{p}_j\right)\\
        & = -(\zeta - 1)\sum_{i=1}^N \lambda_i  \hat{A}_i +\sum_{i=1}^N(\zeta - 1)  \left(\sum_{j=1}^N\frac{p_jX_{ij}}{pY} \hat{p}_j\right) &\text{(Inverting the two sums)} \\
        & = -(\zeta - 1)\sum_{i=1}^N \lambda_i  \hat{A}_i +\sum_{j=1}^N(\zeta - 1)  \hat{p}_j\frac{p_j}{pY} \left(\sum_{i=1}^N X_{ij}\right)\\
        & = -(\zeta - 1)\sum_{i=1}^N \lambda_i  \hat{A}_i +\sum_{j=1}^N(\zeta - 1)  \hat{p}_j \frac{p_j}{pY} X_j\\
        & = -(\zeta - 1)\sum_{i=1}^N \lambda_i  \hat{A}_i +\sum_{j=1}^N(\zeta - 1)  \lambda_j \hat{p}_j\\
           & = -(\zeta - 1)\sum_{i=1}^N \lambda_i  \hat{A}_i +\sum_{i=1}^N(\zeta - 1)  \lambda_i \hat{p}_i\\
         & = -(\zeta - 1)\sum_{i=1}^N \lambda_i  \hat{A}_i +\sum_{i=1}^N \lambda_i \hat{\tilde{A}}_i\\
            & = (-\zeta + 2)\hat{A}\\
        \Rightarrow & \hat{A}(1+\zeta-2)=0\\
        \Rightarrow & \zeta=1.
\end{align*}

\textcolor{orange}{
\textbf{Hence, what we have proven here} is that in order for the definition of industry level productivity to be equal, with the particular assumptions  $\forall i,~\hat{X}_i=0$ and $\forall f, ~ \hat{w}_f =0$, it is necessary to have $\zeta=1$.}

When we replace $\hat{p}_i$ with its dual definition in line 4 of the previous equation, we implicitly write
\begin{align*}
   & \hat{\tilde{A}}_i=(\zeta-1)\hat{p}_i=(\zeta-1)\left(-\hat{A}_i + \sum_{j=1}^N  a_{ji} \hat{p}_j\right)=(\zeta-1)\left(-\hat{\tilde{A}}_i + \sum_{j=1}^N  a_{ji} \hat{p}_j\right)\\
   \Rightarrow &  \hat{\tilde{A}}_i(1+(\zeta-1))=(\zeta-1)\sum_{j=1}^N  a_{ji} \hat{p}_j\\
   \Rightarrow &  \zeta\hat{\tilde{A}}_i=(\zeta-1)\sum_{j=1}^N  a_{ji} \hat{p}_j\\
   \Rightarrow &  \hat{\tilde{A}}_i=\frac{\zeta-1}{\zeta}\sum_{j=1}^N  a_{ji} \hat{p}_j=\hat{A}_i.
\end{align*}

We can rewrite it 
\begin{equation*}
    \frac{\zeta-1}{\zeta}\hat{\Sigma}_{p}=\hat{\Sigma}_{X}=0.
\end{equation*}
Hence $\hat{\tilde{A}}_i=0$. 

We reach the same result since $\zeta=1$, then
\begin{equation*}
    \forall i, \frac{\zeta-1}{\zeta}\sum_{j=1}^N  a_{ji} \hat{p}_j=0=\hat{\tilde{A}}_i=\hat{A}_i.
\end{equation*}

\textcolor{orange}{
\textbf{So what we are saying here} is that in order for the assumption that $\zeta$ is the same across all sector to hold we can only talk about the case where there is no productivity growth.}


% Is it a tautology? Should we be using the dual definition of industry level productivity with \textit{at the same time}, the assumption that $\zeta$ is the same across all sectors (which allows us to take it out of the sum, thus to arrive at 0?) Maybe that's the problem. 

% But isn't also what we are doing when computing downstream and upstream productivity? 


\subsubsection{Without extra assumptions} We have N sectors in the economy. \textbf{No more simplifying assumptions}: $\hat{X}_i\neq 0$ and $ \hat{w}_f \neq 0$.
Hence, industry-level productivity prediction is $\hat{\tilde{A}}_i = \zeta\hat{X}_i+(\zeta - 1) \hat{p}_i$.
Let us concatenate the two equations, assuming that $\hat{A}_i\approx \hat{\tilde{A}}_i$ (which we prove numerically). Let us remember that the dual definition of productivity gives $\hat{A}_i \equiv  - \left( \hat{p}_i - \sum_{j=1}^N  a_{ji} \hat{p}_j - \sum_{f=1}^F \tilde{l}_{fi} \hat{w}_f \right)$.
Hence,
\begin{equation*}
     \hat{p}_i=-\hat{A}_i + \sum_{j=1}^N  a_{ji} \hat{p}_j+\sum_{f=1}^F \tilde{l}_{fi} \hat{w}_f.
\end{equation*}
And the primal, $\hat{A}_i \equiv \hat{X}_i - \sum_{j=1}^N  a_{ji} \hat{X}_{ij} - \sum_{f=1}^F \tilde{l}_{fi}\hat{L}_{fi}$, then
\begin{equation*}
     \hat{X}_i=\hat{A}_i + \sum_{j=1}^N  a_{ji} \hat{X}_{ij} + \sum_{f=1}^F \tilde{l}_{fi}\hat{L}_{fi}.
\end{equation*}

\begin{align*}
    \hat{A}&=\sum_{i=1}^N \lambda_i \hat{A}_i\\
    & \approx \sum_{i=1}^N \lambda_i \hat{\tilde{A}}_i\\
    & = \sum_{i=1}^N \lambda_i \left(\zeta\hat{X}_i+(\zeta - 1) \hat{p}_i\right) \\
     & = \zeta\sum_{i=1}^N \lambda_i\hat{X}_i+(\zeta - 1) \sum_{i=1}^N \lambda_i\hat{p}_i \\
      & = \zeta\sum_{i=1}^N \lambda_i \left(\hat{A}_i + \sum_{j=1}^N  a_{ji} \hat{X}_{ij} + \sum_{f=1}^F \tilde{l}_{fi}\hat{L}_{fi} \right)+(\zeta - 1) \sum_{i=1}^N \lambda_i\left(-\hat{A}_i + \sum_{j=1}^N  a_{ji} \hat{p}_j+\sum_{f=1}^F \tilde{l}_{fi} \hat{w}_f\right) \\
      & = \sum_{i=1}^N \lambda_i\hat{A}_i
      +\zeta \underbrace{\left(
      \sum_{i=1}^N \sum_{j=1}^N \lambda_i a_{ji} \hat{X}_{ij} + \sum_{i=1}^N\sum_{f=1}^F  \lambda_i\tilde{l}_{fi}\hat{L}_{fi}
      \right)}_{A}+
      (\zeta - 1) \underbrace{\left(
      \sum_{i=1}^N \sum_{j=1}^N   \lambda_ia_{ji} \hat{p}_j+\sum_{i=1}^N \sum_{f=1}^F  \lambda_i\tilde{l}_{fi} \hat{w}_f
      \right) }_{B}
\end{align*}

We simplify this equation, using the following
\begin{equation}
    \begin{cases}
\hat{A} \equiv \sum_{i=1}^N \frac{p_iC_i}{pY} \hat{C}_i - \sum_{f=1}^F \frac{w_f L_f}{pY} \hat{L}_f,\\
 L_f \hat{L}_f = \sum_i  L_{fi} \hat{L}_{fi},\\
    X_j=\sum_i X_{ij} + C_j,\\
    \sum_i X_{ij}\hat{X}_{ij} = X_j (\hat{X}_j - (C_j/X_j) \hat{C}_j).
    \end{cases}
\end{equation}



We can simplify term A,
\begin{align*}
\left(
      \sum_{i=1}^N \sum_{j=1}^N \lambda_i a_{ji} \hat{X}_{ij} + \sum_{i=1}^N\sum_{f=1}^F  \lambda_i\tilde{l}_{fi}\hat{L}_{fi}
      \right)&=\left(
      \sum_{i=1}^N \sum_{j=1}^N \frac{p_iX_i}{pY} \frac{p_jX_{ij}}{p_iX_i} \hat{X}_{ij} + \sum_{i=1}^N\sum_{f=1}^F  \frac{p_iX_i}{pY} \frac{w_fL_{fi}}{p_iX_i}\hat{L}_{fi}
      \right)\\
      &=
     \sum_{i=1}^N \sum_{j=1}^N \frac{p_jX_{ij}}{pY} \hat{X}_{ij} + \sum_{i=1}^N\sum_{f=1}^F  \frac{w_fL_{fi}}{pY}\hat{L}_{fi}\\
     &=
      \sum_{j=1}^N\frac{p_j}{pY}\sum_{i=1}^N   X_{ij}\hat{X}_{ij} +\sum_{f=1}^F  \frac{w_f}{pY}  \sum_{i=1}^N L_{fi}\hat{L}_{fi}\\
      &=
      \sum_{j=1}^N\frac{p_j}{pY}\left(X_j (\hat{X}_j - (C_j/X_j) \hat{C}_j)\right) +\sum_{f=1}^F  \frac{w_f}{pY}   L_f \hat{L}_f\\
      &=
      \sum_{j=1}^N\lambda_j\hat{X}_j - \sum_{j=1}^N\frac{p_jC_j}{pY}  \hat{C}_j +\sum_{f=1}^F  \frac{w_f L_f }{pY}  \hat{L}_f.\\
\end{align*}


We can simplify term B,
\begin{align*} 
    \left(
      \sum_{i=1}^N \sum_{j=1}^N  \frac{p_iX_i}{pY} \frac{p_jX_{ij}}{p_iX_i} \hat{p}_j+\sum_{i=1}^N \sum_{f=1}^F  \frac{p_iX_i}{pY} \frac{w_fL_{fi}}{p_iX_i} \hat{w}_f
      \right)
      &=\sum_{i=1}^N \sum_{j=1}^N   \frac{p_jX_{ij}}{pY} \hat{p}_j+\sum_{i=1}^N \sum_{f=1}^F \frac{w_fL_{fi}}{pY} \hat{w}_f\\
      &=
       \sum_{j=1}^N\hat{p}_j \frac{p_j}{pY}\sum_{i=1}^N   X_{ij} + \sum_{f=1}^F \hat{w}_f\frac{w_f}{pY}\sum_{i=1}^N L_{fi} \\
       &=
       \sum_{j=1}^N\hat{p}_j \frac{p_j}{pY}X_{j} + \sum_{f=1}^F \hat{w}_f\frac{w_f}{pY}L_{f} \\
        &=
       \sum_{j=1}^N\lambda_j \hat{p}_j + \sum_{f=1}^F \frac{w_fL_{f}}{pY}\hat{w}_f \\
\end{align*}


Combining all equations back together, 

\begin{align*}
    \hat{A}&  = \sum_{i=1}^N \lambda_i\hat{A}_i
      +\zeta \left(
      \sum_{j=1}^N\lambda_j\hat{X}_j - \sum_{j=1}^N\frac{p_jC_j}{pY}  \hat{C}_j +\sum_{f=1}^F  \frac{w_f L_f }{pY}  \hat{L}_f
      \right)+
      (\zeta - 1) \left(
       \sum_{j=1}^N\hat{p}_j \lambda_j + \sum_{f=1}^F \frac{w_fL_{f}}{pY}\hat{w}_f
      \right) \\
       & = \sum_{i=1}^N \lambda_i\hat{A}_i
      + \sum_{j=1}^N\lambda_j(\zeta \hat{X}_j+ (\zeta - 1)\hat{p}_j)
      +\zeta \left(
     - \sum_{j=1}^N\frac{p_jC_j}{pY}  \hat{C}_j +\sum_{f=1}^F  \frac{w_f L_f }{pY}  \hat{L}_f
      \right)+
      (\zeta - 1) \left(
       \sum_{f=1}^F \frac{w_fL_{f}}{pY}\hat{w}_f
      \right) \\
       & = \sum_{i=1}^N \lambda_i\hat{A}_i
      + \sum_{j=1}^N\lambda_j\hat{\tilde{A}}_j
      +\zeta \left(
    -\hat{A}
      \right)+
      (\zeta - 1) \left(
       \sum_{f=1}^F \frac{w_fL_{f}}{pY}\hat{w}_f
      \right) \\
      \Rightarrow & \hat{A}(1-2+\zeta)= 
      (\zeta - 1) \left(
       \sum_{f=1}^F \frac{w_fL_{f}}{pY}\hat{w}_f
      \right) \\
      \Rightarrow & \hat{A}(\zeta-1)= 
      (\zeta - 1) \left(
       \sum_{f=1}^F \frac{w_fL_{f}}{pY}\hat{w}_f
      \right).
\end{align*}
\textcolor{orange}{Hence,
either
\begin{equation}
    \begin{cases}
        \zeta=1\\
        \textbf{ or } \\
        \hat{A}=  \sum_{f=1}^F \frac{w_fL_{f}}{pY}\hat{w}_f.
    \end{cases}
\end{equation}}


\subsubsection{Conclusion} 
We have not made much progress. The assumption that the two expressions for aggregate productivity yield the same result with a universal $\zeta$ implies one of two possibilities. Either $\zeta = 1$, which means that sales growth is driven entirely by price changes, with no contribution from quantity variations (equivalent to assuming $\forall i,~\hat{X}_i=0$), or aggregate productivity ultimately depends only on changes in factor prices, which are reflected in sectoral prices. 

In other words, in the second case, if we know how wages and the cost of capital evolve, along with their respective shares in GDP, can we fully determine aggregate productivity growth? 

Given this, I am uncertain about the interpretation of the results if we were to test the following equation on historical data:
\begin{equation*}
    \hat{A}=\sum_{i=1}^N\lambda_i (\zeta\hat{X}_i+(\zeta-1)\hat{p}_i).
\end{equation*}


\newpage

\subsection{Solution \#2}

\ER{The main difference with the previous section is that here $\hat{\tilde{A}}_i\neq\hat{A}_i$ because $\zeta_i$ is industry-specific. In this case: which solution shall we prefer ?}

\ER{Note that saying that we would have calibrated $\zeta$ on a more general cases at the aggregate level would not give us an obvious answer, as the shocks are not comparable and we have insider knowledge on the net zero transition (see next section Solution \#3).
}



\ER{First we discuss the case where only an upstream sector undergoes productivity changes (can be made general case by assuming M is all the sectors that are not U aggregated (since everything is linear, it works well).
Second, we consider the case where both upstream and downstream productivity growth. }

\paragraph{Only upstream TFP changes\\}
With a toy model with two sectors (that we developed in another document, refer to it for more details.), downstream and \textbf{\textit{upstream}, with only the \textit{latter} undergoing productivity changes, }we can get the key insights (and not full solution) that if we consider the price change in the downstream sector according to our framework 
\begin{equation}
    \hat{\tilde{A}}_M =  \zeta \hat{X}_M + (\zeta - 1) \hat{p}_M.
\end{equation}
Then the aggregate productivity is
\begin{equation}
      \hat{A} 
      \equiv 
    \lambda_U  \hat{A}_U  +  \lambda_M  \hat{A}_M
     =  \lambda_U  \hat{A}_U=  \lambda_M  \hat{\tilde{A}}_M =\hat{A}.
\end{equation}
Meaning that if we consider only the downstream industry price evolution, we capture the whole aggregate productivity. 


\paragraph{What if we have no prior knowledge of which sector to use}
If we use the price increase and quantities increase of both U and M, then we compute
\begin{align*}
 \hat{\tilde{A}}=\lambda_M  \hat{\tilde{A}}_M +  \lambda_U  \hat{\tilde{A}}_U 
    &=
    \lambda_M  \hat{A}_M +  2\lambda_U  \hat{A}_U.
\end{align*}

If we only consider U, then we predict
\begin{equation*}
 \hat{\tilde{A}}=  \lambda_U  \hat{\tilde{A}}_U= \lambda_U  \hat{A}_U.
\end{equation*}

Let us assume we consider either U, M or both U and M to apply how framework to compute our prediction of aggregate productivity growth $\hat{\tilde{A}}$. In which case are we making the biggest error $\varepsilon=\hat{A}-\hat{\tilde{A}}$?

\begin{equation}
\begin{cases}
    U:~ \hat{\tilde{A}}=  \lambda_U  \hat{A}_U, ~\varepsilon=0,\\
      M:~ \hat{\tilde{A}}=  \lambda_U \hat{A}_U, ~\varepsilon= 0,\\
        U\&M:~ \hat{\tilde{A}}=   2\lambda_U  \hat{A}_U , ~\varepsilon= \lambda_U  \hat{A}_U.\\
\end{cases}
\end{equation}

If we consider the two sectors, then we are making a relative error on aggregate productivity 
\begin{equation*}
   \frac{\varepsilon}{\hat{A}}= \frac{\lambda_U  \hat{A}_U}{ 2\lambda_U \hat{A}_U}=\frac{1}{2}.
\end{equation*}

We are overestimating the aggregate impact of the the U sector on productivity by 50\%. Which is big, but since it is likely that $\lambda_U<<1$, then one may consider that it is under control.


\vspace{3mm}
\paragraph{Both Upstream and Downstream TFP changes\\}

Unlike in the previous section, let us consider that \textbf{there is some idiosyncratic productivity growth in M}.

Consider the simplified economy described, such as
\begin{itemize}
    \item \textbf{Upstream sector $U$} using only labour. 
    \[
    \begin{cases}
    \hat{X}_U > 0,\\
    \hat{p}_U < 0,\\
    0 < \zeta < 1.
    \end{cases}
    \]
    \item \textbf{Downstream sector $M$} using $U$ as its sole input. Let $a_{UM}$ be the cost share of $U$ in $M$'s expenses.  By using primal and dual expressions in productivity growth accounting we can express
    \[
    \begin{cases}
    \hat{X}_M = \hat{A}_M+a_{UM}  \hat{X}_U,\\
    \hat{p}_M = -\hat{A}_M+a_{UM}  \hat{p}_U,\\
    \hat{A}_M > 0.
    \end{cases}
    \]
\end{itemize}

\paragraph{Relating Cost Shares and Domar Weights}
The share of $U$ in $M$'s total expenses is
\begin{equation*}
     a_{UM}  =  \frac{p_U X_U}{p_M X_M}  =  \frac{p_U X_U}{p_Y Y} \frac{p_Y Y}{p_M X_M} 
     =  \frac{\lambda_U}{\lambda_M},
\end{equation*}
where $\lambda_U = \tfrac{p_U X_U}{Y}$ and $\lambda_M = \tfrac{p_M X_M}{Y}$ are the Domar weights of $U$ and $M$, with $p_Y Y$ is total nominal GDP of the economy. Hence, 
\begin{equation*}
    \hat{X}_M = \hat{A}_M+\frac{\lambda_U}{\lambda_M} \hat{X}_U
    \quad\text{and}\quad
    \hat{p}_M = -\hat{A}_M+\frac{\lambda_U}{\lambda_M} \hat{p}_U.
\end{equation*}

\paragraph{True Aggregate Productivity.}
Aggregate productivity growth $\hat{A}$ is 
\begin{equation}
      \hat{A} 
     =  
    \lambda_U  \hat{A}_U  +  \lambda_M  \hat{A}_M.
\end{equation}


\paragraph{Downstream-Only Productivity Measure} We now show that a simplified measure of productivity that focuses solely on the downstream sector $M$ still recovers $\hat{A}$. In our framework, we define a downstream ``effective'' productivity growth $\hat{\tilde{A}}_M$ as
\begin{equation*}
    \hat{\tilde{A}}_M =  \zeta \hat{X}_M + (\zeta - 1) \hat{p}_M.
\end{equation*}
Then the overall productivity, from the perspective of $M$ alone, is scaled by the Domar weight $\lambda_M$:
\begin{equation}
    \begin{aligned}
    \hat{\tilde{A}}=\lambda_M  \hat{\tilde{A}}_M 
    &= \lambda_M\Bigl(\zeta \hat{X}_M + (\zeta-1) \hat{p}_M\Bigr)\\
    &= \lambda_M\Bigl(\zeta(\hat{A}_M+ \tfrac{\lambda_U}{\lambda_M} \hat{X}_U )
        +  (\zeta-1) (-\hat{A}_M+\tfrac{\lambda_U}{\lambda_M} \hat{p}_U)\Bigr)\\
    &= \lambda_M (\zeta\hat{A}_M-(\zeta-1)\hat{A}_M+)+\lambda_U\Bigl(\zeta \hat{X}_U + (\zeta-1) \hat{p}_U\Bigr)\\
    &= \lambda_M \hat{A}_M+\lambda_U \hat{\tilde{A}}_U\\
    &=\lambda_M \hat{A}_M+\lambda_U \hat{A}_U\quad(\text{since we calibrate }\zeta\text{ such that }\hat{\tilde{A}}_U=\hat{A}_U).
\end{aligned}
\end{equation}
We know that $\hat{\tilde{A}}_U = \hat{A}_U$ because we calibrate $\zeta$ to exactly match observed productivity growth.
This confirms that our method correctly captures the full productivity dynamics by considering only the downstream industry. The reason it works relies on two key points: (i) in standard growth accounting, sector $M$'s productivity growth does not reflect the price decline of sector $U$, whereas here we account for it as part of $M$'s productivity; (ii) using $\lambda_M a_{UM} = \lambda_U$, we ensure the correct Domar weight appears in the decomposition. 

Hence,
\begin{equation}
      \hat{A}  =  \lambda_U \hat{A}_U+\lambda_M \hat{A}_M 
     =  \lambda_M \hat{\tilde{A}}_M.
\end{equation}
In short, analysing the downstream sector alone (with a suitable measure of productivity growth) captures the true aggregate productivity growth.

\paragraph{What if we use the two sectors?}
If we use the price increase and quantities increase of both U and M, then we compute
\begin{align*}
 \hat{\tilde{A}}=\lambda_M  \hat{\tilde{A}}_M +  \lambda_U  \hat{\tilde{A}}_U 
    &=
    \lambda_M  \hat{A}_M +  2\lambda_U  \hat{A}_U.
\end{align*}

If we only consider U, then we predict
\begin{equation*}
 \hat{\tilde{A}}=  \lambda_U  \hat{\tilde{A}}_U= \lambda_U  \hat{A}_U.
\end{equation*}

Let us assume we consider either U, M or both U and M to apply how framework to compute our prediction of aggregate productivity growth $\hat{\tilde{A}}$. In which case are we making the biggest error $\varepsilon=\hat{A}-\hat{\tilde{A}}$?

\begin{equation}
\begin{cases}
    U:~ \hat{A} \approx  \lambda_U  \hat{A}_U, ~\varepsilon= -\lambda_M  \hat{A}_M,\\
      M:~ \hat{A} \approx  \lambda_U \hat{A}_U+\lambda_M \hat{A}_M , ~\varepsilon= 0,\\
        U\&M:~ \hat{A} \approx  \lambda_M  \hat{A}_M +  2\lambda_U  \hat{A}_U , ~\varepsilon= \lambda_U  \hat{A}_U.\\
\end{cases}
\end{equation}

If we consider the two sectors, then we are making a relative error on aggregate productivity 
\begin{equation*}
   \frac{\varepsilon}{\hat{A}}= \frac{\lambda_U  \hat{A}_U}{ \lambda_U \hat{A}_U+\lambda_M \hat{A}_M}<<1
\end{equation*}
which is small if $\lambda_U<<\lambda_M$, which is likely to be the case and unlikely to be compensated by the magnitude of the relative productivity shock.

Hence, the conclusion is, it's better to consider the ``downstream'' sector which has both input changes and idiosyncratic productivity shock if we have prior information on it. If not, then the error we're making is still relatively small. 

\paragraph{Conclusion \# 1}\er{Yes, there is a risk of double counting productivity growth if we consider both U and M. But the error is under control if one sector is smaller than the other.} 


\subsection{Solution \# 3} \textbf{Two part answer. }


\paragraph{In the historical data}. First, $\zeta\approx0.13$ means that's there is strong correlation between each industry's price changes and its productivity growth. In the case where for a majority of industries their price changes would only reflect input price changes, then this correlation would be much lower. Second, on historical data, we are considering trends rather than shocks. Meaning that during the period that we are considering (1980s-2020s), there were multiple innovations and evolutions for each individual sector, meaning that at it is likely that the prices changes of each aggregated industry would reflect some idiosyncratic productivity growth. Third, we stress that the calibration of $\zeta$ takes place only at the industry level. There is no control for the aggregate productivity. Fourth, we then need to sum up all sectors to test for Domar aggregation because \textit{all} sectors have idiosyncratic productivity growth.

\paragraph{In the case of the Net Zero} In Net zero scenarios, we are considering a \textit{shock} (and not a trend) that does not affect all sectors (both extensive (some sectors are not impacted) and intensive (not all sectors are impacted the same) margins). Hence, the situation is a bit different than what we have been estimating (and not sure we can do better, except if we can isolate a single shock similar to the net zero transition, but that is one of our limits in the slides ``no past data on such a transition''). 

To use the framework as we have estimated it: 
\begin{enumerate}
    \item We need to consider aggregated industries big industries ($\zeta$ is calibrated on M rather on U).
    \item If we know (prior on the transition) that one sector is not going to be more productivity but will only reflect input changes (e.g. services reflecting changes in energy prices), then we ought not to consider them (to minimise the error). 
    \item If a big sector is likely to undergo both idiosyncratic productivity shock and input price shocks, then the clever way to handle it is to consider only this big sector (M here)
\end{enumerate}

Meaning that if know which sectors are likely to be affected through their inputs or through their own productivity, then it makes sense to adapt our framework to it. If we don't then we bias our productivity estimates. \ER{The remaining question is: is this bias' magnitude significant? I don't think so because the ``downstream'' sector is typically larger. It comes back to the strength of the Domar aggregation: we can disregard small sectors.}


\section{Appendix: Why are we only considering the downstream industry}
\label{appendix_downstream_industry}

In net-zero transition analyses, certain sectors (e.g., energy) are considered ``impacted'' by policy and technology shifts. However, should the analysis also include upstream industries (e.g., PV, battery, or heat pump manufacturing) rather than focusing solely on downstream sectors?

We focus on the downstream sector to avoid double counting productivity gains. If sector $M$ benefits from productivity improvements in upstream sector $U$, the cost of expanding $U$ is offset by increased sales to $M$. Treating upstream costs separately would double count investments already financed by downstream expenditures. Since our framework tracks only output price growth, including $U$ directly would risk double counting its productivity gains in both $U$ and $M$. It is similar to the processing of CCC extra transition cost to remove double counting. Since our framework focuses only on output price growth, we would be at risk of double counting the productivity gains of $U$ in both $U$ and $M$.

The rest of this appendix formalises why considering the downstream sector alone suffices to capture overall productivity shifts from upstream price or technology changes, extending the argument to multiple downstream industries.

\subsection{Key Insight}
\label{appendix_onepager}


Let the economy be composed of two sectors:
\begin{itemize}
    \item An upstream industry $U$ (e.g., producing solar PV modules, batteries).
    \item A downstream sector $M$ (e.g., the electricity generation sector that purchases solar PV from $U$).
\end{itemize}
\paragraph{Core Question} If $M$ appears more productive (i.e. offers a lower price) only because it purchases cheaper inputs from $U$, do we need to model $U$ explicitly to capture the net productivity gains? Or is it enough to consider $M$ only?

\paragraph{Short Answer} 
It is enough to consider the downstream sector, because the productivity gains of the upstream sectors are accurately reflected in the downstream price dynamics. 

\paragraph{Two-Sector economy} Consider the simplified economy described, such as
\begin{itemize}
    \item \textbf{Upstream sector $U$} using only labour. 
    \[
    \begin{cases}
    \hat{X}_U > 0,\\
    \hat{p}_U < 0,\\
    0 < \zeta < 1.
    \end{cases}
    \]
    \item \textbf{Downstream sector $M$} using $U$ as its sole input. Let $a_{UM}$ be the cost share of $U$ in $M$'s expenses. By assumption, there is no idiosyncratic productivity growth in M. By using primal and dual expressions in productivity growth accounting we can express
    \[
    \begin{cases}
    \hat{X}_M = a_{UM}  \hat{X}_U,\\
    \hat{p}_M = a_{UM}  \hat{p}_U,\\
    \hat{A}_M = 0.
    \end{cases}
    \]
\end{itemize}

\paragraph{Relating Cost Shares and Domar Weights}
The share of $U$ in $M$'s total expenses is
\begin{equation*}
     a_{UM}  =  \frac{p_U X_U}{p_M X_M}  =  \frac{p_U X_U}{p_Y Y} \frac{p_Y Y}{p_M X_M} 
     =  \frac{\lambda_U}{\lambda_M},
\end{equation*}
where $\lambda_U = \tfrac{p_U X_U}{Y}$ and $\lambda_M = \tfrac{p_M X_M}{Y}$ are the Domar weights of $U$ and $M$, with $p_Y Y$ is total nominal GDP of the economy. Hence, 
\begin{equation*}
    \hat{X}_M = \frac{\lambda_U}{\lambda_M} \hat{X}_U
    \quad\text{and}\quad
    \hat{p}_M = \frac{\lambda_U}{\lambda_M} \hat{p}_U.
\end{equation*}

\paragraph{True Aggregate Productivity.}
Aggregate productivity growth $\hat{A}$ is 
\begin{equation}
      \hat{A} 
     =  
    \lambda_U  \hat{A}_U  +  \lambda_M  \hat{A}_M
     =  \lambda_U  \hat{A}_U
    \quad (\text{since }\hat{A}_M=0).
\end{equation}


\paragraph{Downstream-Only Productivity Measure} We now show that a simplified measure of productivity that focuses solely on the downstream sector $M$ still recovers $\hat{A}$. In our framework, we define a downstream ``effective'' productivity growth $\hat{\tilde{A}}_M$ as
\begin{equation*}
    \hat{\tilde{A}}_M =  \zeta \hat{X}_M + (\zeta - 1) \hat{p}_M.
\end{equation*}
Then the overall productivity, from the perspective of $M$ alone, is scaled by the Domar weight $\lambda_M$:
\begin{equation}
    \begin{aligned}
    \lambda_M  \hat{\tilde{A}}_M 
    &= \lambda_M\Bigl(\zeta \hat{X}_M + (\zeta-1) \hat{p}_M\Bigr)\\
    &= \lambda_M\Bigl(\zeta \tfrac{\lambda_U}{\lambda_M} \hat{X}_U 
        +  (\zeta-1) \tfrac{\lambda_U}{\lambda_M} \hat{p}_U\Bigr)\\
    &= \lambda_U\Bigl(\zeta \hat{X}_U + (\zeta-1) \hat{p}_U\Bigr)\\
    &= \lambda_U \hat{\tilde{A}}_U\\
    &=\lambda_U \hat{A}_U\quad(\text{since we calibrate }\zeta\text{ such that }\hat{\tilde{A}}_U=\hat{A}_U).
\end{aligned}
\end{equation}
We know that $\hat{\tilde{A}}_U = \hat{A}_U$ because we calibrate $\zeta$ to exactly match observed productivity growth.
This confirms that our method correctly captures the full productivity dynamics by considering only the downstream industry. The reason it works relies on two key points: (i) in standard growth accounting, sector $M$'s productivity growth does not reflect the price decline of sector $U$, whereas here we account for it as part of $M$'s productivity; (ii) using $\lambda_M a_{UM} = \lambda_U$, we ensure the correct Domar weight appears in the decomposition. 

Hence,
\begin{equation}
      \hat{A}  =  \lambda_U \hat{A}_U  =  \lambda_U \hat{\tilde{A}}_U 
     =  \lambda_M \hat{\tilde{A}}_M.
\end{equation}
In short, analysing the downstream sector alone (with a suitable measure of productivity growth) captures the true aggregate productivity growth.

\FL{(We probably need to clarify the argument. I wouldn't say our equation gives a ``measure'' of productvity growth; more like a ``prediction''. )}

\FL{What happens if we are given price and output growth for the U sector? Then our frameowrk gives a biased aggregate productivity (exactly double, right?). So we should be careful to write clearly the conclusion that our framework risks double counting! Does this not contradict our result that empirically, we use productivity shocks from all industries in KLEMS to get unbiased results?}

\subsection{Discussion: Why Not Model the Upstream Sector Directly?}
Downstream sectors, such as energy or transport, often benefit from multiple upstream productivity gains (e.g., solar PV, battery manufacturing, wind turbines). Tracking each upstream sector separately would require detailed knowledge of cost shares and demand flows. The streamlined method shown here implicitly accounts for changes in upstream prices within the downstream productivity and price data. 


\subsection{Extensions to Multiple Downstream Sectors}
Consider now an upstream sector $U$ supplying two downstream sectors $M$ and $C$. Suppose $U$ uses only labour and sells to $M$ and $C$ in proportions $\beta_M$ and $\beta_C$, such as $\beta_M + \beta_C = 1$. We denote
\[
\begin{aligned}
    &X_{MU} = \beta_M X_U, 
    \quad 
    X_{CU} = \beta_C X_U, 
    \quad \text{and}  \quad 
    X_U = X_{MU} + X_{CU}.
\end{aligned}
\]
If both downstream sectors use $U$ as their only input,\footnote{
%
This can be generalized to situations where $M$ and $C$ have other inputs. The same logic holds so long as $\zeta$ is appropriately calibrated. \FL{What if the network has loops?}
%
} we have
\[
\begin{cases}
\hat{X}_U>0,\\
\hat{p}_U<0,\\
0<\zeta<1,
\end{cases}
\quad
\begin{cases}
\hat{X}_M&=a_{UM}\hat{X}_U,\\
\hat{p}_M&=a_{UM}\hat{p}_U,\\
\hat{A}_M&=0,
\end{cases}
\quad
\begin{cases}
\hat{X}_C&=a_{UC}\hat{X}_U,\\
\hat{p}_C&=a_{UC}\hat{p}_U,\\
\hat{A}_C&=0.
\end{cases}
\]

Here, 
\begin{equation*}
    a_{UM}=\frac{X_{UM}p_U}{p_MX_M}=\frac{\beta_CX_Up_U}{p_MX_M}=\beta_C\lambda_U
\end{equation*}
Thus aggregate productivity growth is
\begin{equation}
\begin{aligned}
    \hat{A}&=\lambda_M \hat{\tilde{A}}_M+\lambda_C \hat{\tilde{A}}_C\\
    &= \lambda_M \left(\zeta\hat{X}_M+(\zeta-1)\hat{p}_M\right)+
    \lambda_C \left(\zeta\hat{X}_C+(\zeta-1)\hat{p}_C\right)\\
    &=\lambda_M \left(\zeta\hat{X}_U\beta_M\frac{\lambda_U}{\lambda_M}+(\zeta-1)\hat{p}_U\beta_M\frac{\lambda_U}{\lambda_M}\right)+
    \lambda_C \left(\zeta\hat{X}_U\beta_C\frac{\lambda_U}{\lambda_C}+(\zeta-1)\hat{p}_U\beta_C\frac{\lambda_U}{\lambda_C}\right)\\
     &=\lambda_U(\beta_M+\beta_C) \left(\zeta\hat{X}_U+(\zeta-1)\hat{p}_U\right)\\
     &=\lambda_U \hat{\tilde{A}}_U,\\
\end{aligned}
\end{equation}

Thus, tracking price and production trends across all downstream sectors that depend on $U$ is sufficient to fully capture the upstream sector's contribution to productivity. More generally, we should account for all downstream sectors affected by an upstream sector, provided that the input is significant enough to meaningfully influence downstream prices and quantities.


\subsection{Variation on the two sector model}

\ER{New section}
\paragraph{Two-Sector economy} Consider the simplified economy described, such as
\begin{itemize}
    \item \textbf{Upstream sector $U$} using only labour. 
    \[
    \begin{cases}
    \hat{X}_U > 0,\\
    \hat{p}_U < 0,\\
    0 < \zeta < 1.
    \end{cases}
    \]
    \item \textbf{Downstream sector $M$} using $U$ as its sole input. Let $a_{UM}$ be the cost share of $U$ in $M$'s expenses. Unlike in the previous section, there is some idiosyncratic productivity growth in M. By using primal and dual expressions in productivity growth accounting we can express
    \[
    \begin{cases}
    \hat{X}_M = \hat{A}_M+a_{UM}  \hat{X}_U,\\
    \hat{p}_M = -\hat{A}_M+a_{UM}  \hat{p}_U,\\
    \hat{A}_M > 0.
    \end{cases}
    \]
\end{itemize}

\paragraph{Relating Cost Shares and Domar Weights}
The share of $U$ in $M$'s total expenses is
\begin{equation*}
     a_{UM}  =  \frac{p_U X_U}{p_M X_M}  =  \frac{p_U X_U}{p_Y Y} \frac{p_Y Y}{p_M X_M} 
     =  \frac{\lambda_U}{\lambda_M},
\end{equation*}
where $\lambda_U = \tfrac{p_U X_U}{Y}$ and $\lambda_M = \tfrac{p_M X_M}{Y}$ are the Domar weights of $U$ and $M$, with $p_Y Y$ is total nominal GDP of the economy. Hence, 
\begin{equation*}
    \hat{X}_M = \hat{A}_M+\frac{\lambda_U}{\lambda_M} \hat{X}_U
    \quad\text{and}\quad
    \hat{p}_M = -\hat{A}_M+\frac{\lambda_U}{\lambda_M} \hat{p}_U.
\end{equation*}

\paragraph{True Aggregate Productivity.}
Aggregate productivity growth $\hat{A}$ is 
\begin{equation}
      \hat{A} 
     =  
    \lambda_U  \hat{A}_U  +  \lambda_M  \hat{A}_M.
\end{equation}


\paragraph{Downstream-Only Productivity Measure} We now show that a simplified measure of productivity that focuses solely on the downstream sector $M$ still recovers $\hat{A}$. In our framework, we define a downstream ``effective'' productivity growth $\hat{\tilde{A}}_M$ as
\begin{equation*}
    \hat{\tilde{A}}_M =  \zeta \hat{X}_M + (\zeta - 1) \hat{p}_M.
\end{equation*}
Then the overall productivity, from the perspective of $M$ alone, is scaled by the Domar weight $\lambda_M$:
\begin{equation}
    \begin{aligned}
    \hat{\tilde{A}}=\lambda_M  \hat{\tilde{A}}_M 
    &= \lambda_M\Bigl(\zeta \hat{X}_M + (\zeta-1) \hat{p}_M\Bigr)\\
    &= \lambda_M\Bigl(\zeta(\hat{A}_M+ \tfrac{\lambda_U}{\lambda_M} \hat{X}_U )
        +  (\zeta-1) (-\hat{A}_M+\tfrac{\lambda_U}{\lambda_M} \hat{p}_U)\Bigr)\\
    &= \lambda_M (\zeta\hat{A}_M-(\zeta-1)\hat{A}_M+)+\lambda_U\Bigl(\zeta \hat{X}_U + (\zeta-1) \hat{p}_U\Bigr)\\
    &= \lambda_M \hat{A}_M+\lambda_U \hat{\tilde{A}}_U\\
    &=\lambda_M \hat{A}_M+\lambda_U \hat{A}_U\quad(\text{since we calibrate }\zeta\text{ such that }\hat{\tilde{A}}_U=\hat{A}_U).
\end{aligned}
\end{equation}
We know that $\hat{\tilde{A}}_U = \hat{A}_U$ because we calibrate $\zeta$ to exactly match observed productivity growth.
This confirms that our method correctly captures the full productivity dynamics by considering only the downstream industry. The reason it works relies on two key points: (i) in standard growth accounting, sector $M$'s productivity growth does not reflect the price decline of sector $U$, whereas here we account for it as part of $M$'s productivity; (ii) using $\lambda_M a_{UM} = \lambda_U$, we ensure the correct Domar weight appears in the decomposition. 

Hence,
\begin{equation}
      \hat{A}  =  \lambda_U \hat{A}_U+\lambda_M \hat{A}_M 
     =  \lambda_M \hat{\tilde{A}}_M.
\end{equation}
In short, analysing the downstream sector alone (with a suitable measure of productivity growth) captures the true aggregate productivity growth.

\paragraph{What if we use the two sectors?}
If we use the price increase and quantities increase of both U and M, then we compute
\begin{align*}
 \hat{\tilde{A}}=\lambda_M  \hat{\tilde{A}}_M +  \lambda_U  \hat{\tilde{A}}_U 
    &=
    \lambda_M  \hat{A}_M +  2\lambda_U  \hat{A}_U.
\end{align*}

If we only consider U, then we predict
\begin{equation*}
 \hat{\tilde{A}}=  \lambda_U  \hat{\tilde{A}}_U= \lambda_U  \hat{A}_U.
\end{equation*}

Let us assume we consider either U, M or both U and M to apply how framework to compute our prediction of aggregate productivity growth $\hat{\tilde{A}}$. In which case are we making the biggest error $\varepsilon=\hat{A}-\hat{\tilde{A}}$?

\begin{equation}
\begin{cases}
    U:~ \hat{A} \approx  \lambda_U  \hat{A}_U, ~\varepsilon= -\lambda_M  \hat{A}_M,\\
      M:~ \hat{A} \approx  \lambda_U \hat{A}_U+\lambda_M \hat{A}_M , ~\varepsilon= 0,\\
        U\&M:~ \hat{A} \approx  \lambda_M  \hat{A}_M +  2\lambda_U  \hat{A}_U , ~\varepsilon= \lambda_U  \hat{A}_U.\\
\end{cases}
\end{equation}

If we consider the two sectors, then we are making a relative error on aggregate productivity 
\begin{equation*}
   \frac{\varepsilon}{\hat{A}}= \frac{\lambda_U  \hat{A}_U}{ \lambda_U \hat{A}_U+\lambda_M \hat{A}_M}<<1
\end{equation*}
which is small if $\lambda_U<<\lambda_M$, which is likely to be the case and unlikely to be compensated by the magnitude of the relative productivity shock. 

Hence, the conclusion is, it's better to consider the ``downstream'' sector which has both input changes and idiosyncratic productivity shock if we have prior information on it. If not, then the error we're making is still relatively small. 



\subsection{Particular Case: Upstream and Downstream with Fixed Quantities}
\label{sec:particularcase}

\ER{Not sure we should include this section, but footnote 2 is interesting and I would like your opinion François.}
Consider now an extreme scenario:
\begin{itemize}
    \item \textbf{Upstream sector $U$}: 
    \[
    \begin{cases}
    \hat{X}_U = 0,\\
    \hat{p}_U < 0,\\
    \hat{A}_U = - \hat{p}_U > 0
    \end{cases}
    \]
    \item \textbf{Downstream sector $M$}: 
    \[
    \begin{cases}
    \hat{X}_M=0,\\
    \hat{p}_M = - \hat{A}_M  +  a_{UM} \hat{p}_U,\\
    \hat{A}_M = -\bigl(\hat{p}_M - a_{UM} \hat{p}_U\bigr).
    \end{cases}
    \]
\end{itemize}
The aggregate productivity is
\begin{equation}
\begin{aligned}
    \hat{A}&=\lambda_U\hat{A}_U+\lambda_M\hat{A}_M\\
    &=-\frac{p_UX_U}{Y}\hat{p}_U+\frac{p_MX_M}{Y}(-\hat{p}_M + a_{UM}\hat{p}_U)\\
    &=-\frac{p_MX_M}{Y}\hat{p}_U\\
    &=-\lambda_M \hat{p}_M.
\end{aligned}
\end{equation}


On the other hand, our simplified measure would be
\[
    \hat{\tilde{A}}_M 
    = \zeta \hat{X}_M + (\zeta-1) \hat{p}_M.
\]
Here, since $\hat{X}_M = 0$, we calibrate $\zeta=0$ such as  $\hat{\tilde{A}}_M = - \hat{p}_M$. \footnote{
However, from the formal definition 
\[
    \zeta  =  \frac{\hat{\Sigma}_{p_i}}{\hat{\Sigma}_{p_i} + \hat{\Sigma}_{X_i}},
\]

one might instead find that $\zeta=1$ since $\hat{\Sigma}_{X_i}=0$. This discrepancy signals that the parameter $\zeta$ must be carefully calibrated to ensure internal consistency when quantities are fixed. There is a subtle mismatch in this particular edge case.}

\subsection{Conclusion}
From a macroeconomic or aggregate perspective, one may avoid double counting both upstream and downstream costs by (i) tracking downstream prices and production changes and (ii) applying a suitable measure of productivity growth as in our framework. 
\end{document}

\newpage

\section{Appendix: Why are we only considering the downstream industry}
\label{appendix_downstream_industry}

\paragraph{Summary} In our method, we only consider the energy sector, without considering all the upstream industries: batteries, solar PV production, etc. Likewise, for heating we do not consider the manufacturing of heat pumps. 
Two reasons: 
\begin{itemize}
    \item If only the sector M benefits from the productivity gains of an upstream sector U, then it's enough to consider the sector M in isolation. The intutition is the same as for removing the double counting (processing) of the CCC cost. If one manufacturing sector M electrifies its usages, then the extra electricity bills of the sector M compensates the extra investment in new electricity production capacity of the energy sector E. Then, the "extra cost" of the sector E are already financed by extra sales, while the ones of the sector M are actual extra cost of the transition (or vice-versa). 
    \item In the case, where 
\end{itemize}

\ER{How to consider the solar pv production ?}

\subsection{One pager}

\paragraph{Question} Let the economy be composed of two sectors, an upstream industry U -- for instance producing solar PV -- and a downstream sector M -- for instance electricity production sector. A problem is: we are considering some sectors as impacted by the net zero transition, for instance the energy sector, but if it is becoming more productive, is it only because its own merit or because it is buying solar pv and windmills cheaper ? Thus, should we consider the upstream manufacturing sector instead of just the energy sector ? Or both ? 
More generally, should we consider in our approach all the upstream sectors ?

\paragraph{Answer} We can only consider the downstream. 

\paragraph{Proof} 

\begin{itemize}
    \item \textbf{Upstream sector U}: Using only labour. 
    \[
\left\{
\begin{aligned}
\hat{X}_U&>0\\
\hat{p}_U&<0\\
0&<\zeta<1
\end{aligned}
\right.
\]
  \item \textbf{Downstream sector M}:M is only using one input, that is U, Using the dual expression of productivity growth to express $\hat{p}_M$. We note $a_{UM}$ the share of U in M expenses.
    \[
\left\{
\begin{aligned}
\hat{X}_M&=a_{UM}\hat{X}_U\\
\hat{p}_M&=a_{UM}\hat{p}_U\\
\hat{A}_M&=0
\end{aligned}
\right.
\]
\end{itemize}

Prior to looking at the productivity, let us note that the share of U in M's expenses can be expressed as
\begin{align*}
    a_{UM}=\frac{p_UX_U}{p_MX_M}=\frac{p_UX_U}{Y}\frac{Y}{p_MX_M}=\frac{\lambda_U}{\lambda_M}.
\end{align*}
We can then rewrite for the sector M
\begin{equation*}
  \hat{X}_M=  \frac{\lambda_U}{\lambda_M}\hat{X}_U\quad\text{and}\quad
\hat{p}_M=  \frac{\lambda_U}{\lambda_M}\hat{p}_U.
\end{equation*}
The true productivity growth is 
\begin{equation}
    \hat{A}=\lambda_U\hat{A}_U+\lambda_M\hat{A}_M=\lambda_U\hat{A}_U.
\end{equation}

Let us enquire whether our framework is able to identify the true aggregate productivity growth by application to the downstream sector only. In our framework, we only consider the downstream sector using the productivity growth $\hat{\tilde{A}}_M$, which is the productivity growth disregarding the input price growth,
\begin{equation}
\begin{aligned}
    \hat{A}&=\lambda_M \hat{\tilde{A}}_M\\
    &= \lambda_M (\zeta\hat{X}_M+(\zeta-1)\hat{p}_M)\\
    &=\lambda_M (\zeta\hat{X}_U\frac{\lambda_U}{\lambda_M}+(\zeta-1)\hat{p}_U\frac{\lambda_U}{\lambda_M})\\
     &=\lambda_U (\zeta\hat{X}_U+(\zeta-1)\hat{p}_U)\\
     &=\lambda_U \hat{\tilde{A}}_U,\\
\end{aligned}
\end{equation}

And we do know that $\hat{\tilde{A}}_U=\hat{A}_U$ because that how we calibrate the $\zeta$ to match exactly the productivity growth (accounting). It means that our method works, and that by considering only the downstream industry we capture the full productivity dynamics. Why does it work: two tricks,First in pure growth accounting the productivity growth of the sector M is [privée] de la baisse de prix de U. Ici on considère qu'elle participe à la productivité de M, second, $\lambda_M a_{UM}=\lambda_U$, hence we can make appear the right domar weight.

\paragraph{Why not consider directly the upstream sector?} Because the downstream sector: energy, transport, etc, usually benefit from multiple upstream productivity growth and their own, it is simpler to consider altogether the impact on their price and production. If not, it would require to know the cost share of each input, which is precisely what we wanted to escape with this method.

\paragraph{With multiple downstream sectors} here it works with one sector, but with several sectors using U and M and C using multiple inputs, it does work as well. 
Note that we have to consider exact $\zeta$ for the U industry. 


\begin{itemize}
    \item \textbf{Upstream sector U}: Using only labour. Supplying to M and C, in proportion $\beta_M$ and $\beta_C$, such as $\beta_M+\beta_C=1$. 

    $X_{MU}=\beta_M X_U$, $X_{CU}=\beta_C X_U$,
    $X_U=X_{MU}+X_{CU}=X_{CU}\left(\frac{X_{MU}}{X_{CU}}+1\right)=X_{CU}(\gamma+1)=X_{MU}(\frac{1}{\gamma}+1)$
    Hence, $\hat{X}_U=\hat{X}_{CU}=\hat{X}_{MU}$
    \[
\left\{
\begin{aligned}
\hat{X}_U&>0\\
\hat{p}_U&<0\\
0&<\zeta<1
\end{aligned}
\right.
\]
  \item \textbf{Downstream sector M}:M is only using one input, that is U, Using the dual expression of productivity growth to express $\hat{p}_M$. We note $a_{UM}$ the share of U in M expenses.
    \[
\left\{
\begin{aligned}
\hat{X}_M&=a_{UM}\hat{X}_U\\
\hat{p}_M&=a_{UM}\hat{p}_U\\
\hat{A}_M&=0
\end{aligned}
\right.
\]
\item Likewise for C.
   \[
\left\{
\begin{aligned}
\hat{X}_C&=a_{UC}\hat{X}_U\\
\hat{p}_C&=a_{UC}\hat{p}_U\\
\hat{A}_C&=0
\end{aligned}
\right.
\]
\end{itemize}

Here, 
\begin{equation*}
    a_{UM}=\frac{X_{UM}p_U}{p_MX_M}=\frac{\beta_CX_Up_U}{p_MX_M}=\beta_C\lambda_U
\end{equation*}
Thus the aggregate productivity
\begin{equation}
\begin{aligned}
    \hat{A}&=\lambda_M \hat{\tilde{A}}_M+\lambda_C \hat{\tilde{A}}_C\\
    &= \lambda_M (\zeta\hat{X}_M+(\zeta-1)\hat{p}_M)+
    \lambda_C (\zeta\hat{X}_C+(\zeta-1)\hat{p}_C)\\
    &=\lambda_M (\zeta\hat{X}_U\beta_M\frac{\lambda_U}{\lambda_M}+(\zeta-1)\hat{p}_U\beta_M\frac{\lambda_U}{\lambda_M})+
    \lambda_C (\zeta\hat{X}_U\beta_C\frac{\lambda_U}{\lambda_C}+(\zeta-1)\hat{p}_U\beta_C\frac{\lambda_U}{\lambda_C})\\
     &=\lambda_U(\beta_M+\beta_C) (\zeta\hat{X}_U+(\zeta-1)\hat{p}_U)\\
     &=\lambda_U \hat{\tilde{A}}_U,\\
\end{aligned}
\end{equation}

In any case, we should consider all the downstream sectors of an affected upstream sector, at the condition that this input is big enough to have a sizeable impact on the price and quantities of the downstream sector.


\subsubsection{Simple particular case} In this particular case, the upstream industry uses supplies only the manufacturing sector which only supplies household. 
\ER{With one twist, see footnote.}

In a purely accounting framework

\begin{itemize}
    \item \textbf{Upstream sector U}: 
    \[
\left\{
\begin{aligned}
\hat{X}_U&=0\\
\hat{p}_U&<0\\
\hat{A}_U&=-\hat{p}_U>0\\
\end{aligned}
\right.
\]
  \item \textbf{Downstream sector M}: Using the dual expression of productivity growth to express $\hat{p}_M$. $a_{UM}$ is the share of U in M expenses.
    \[
\left\{
\begin{aligned}
\hat{X}_M&=0\\
\hat{p}_M&=-\hat{A}_M + a_{UM}\hat{p}_U\\
\hat{A}_M&=-(\hat{p}_M - a_{UM}\hat{p}_U)
\end{aligned}
\right.
\]
\end{itemize}

Aggregate productivity can be simplified as 
\begin{equation}
\begin{aligned}
    \hat{A}&=\lambda_U\hat{A}_U+\lambda_M\hat{A}_M\\
    &=-\frac{p_UX_U}{Y}\hat{p}_U+\frac{p_MX_M}{Y}(-\hat{p}_M + a_{UM}\hat{p}_U)\\
    &=-\frac{p_MX_M}{Y}\hat{p}_U\\
    &=-\lambda_M \hat{p}_M.
\end{aligned}
\end{equation}

Does it match with our approximation of 
\begin{equation}
    \hat{A}_M=\zeta\hat{X}_M+(\zeta-1)\hat{p}_M?
\end{equation}

Yes, because in this particular economy, since quantities of M are fixed, $\zeta=0$.\footnote{
\ER{Warning point} $\zeta$ needs to be $=0$ to respect the $\hat{\tilde{A}}=-\hat{p}_M$ (with $\hat{\tilde{A}}$ being not the true productivity but the productivity growth computed with our framework). But when looking at the definition of $\zeta$ it's not that obvious. With 
\[
\left\{
\begin{aligned}
\hat{\Sigma}_{p_i}& \equiv \sum_{j=1}^N  w_j \hat{p}_j - \sum_{f=1}^F \tilde{l}_{fi} \hat{w}_f\\
\hat{\Sigma}_{X_i} &\equiv \sum_{j=1}^N  a_{ji} \hat{X}_{ij} - \sum_{f=1}^F  \tilde{l}_{fi}\hat{L}_{fi}\\
\zeta &\equiv \frac{\hat{\Sigma}_{p_i}}{\hat{\Sigma}_{p_i} + \hat{\Sigma}_{X_i}}.
\end{aligned}
\right.
\]
In this economy: 
\[
\left\{
\begin{aligned}
\hat{\Sigma}_{p_i}&=  a_{UM} \hat{p}_U \\
\hat{\Sigma}_{X_i}&=0\\
\zeta &=\frac{\hat{\Sigma}_{p_i}}{\hat{\Sigma}_{p_i}}=1.\\
\end{aligned}
\right.
\]}


% % Bibliography ---------------------------
% \newpage
% \small
% \renewcommand\bibname{References}
% \bibliographystyle{apalike}
% \bibliography{z-bib}
% \normalsize

% \newpage

\end{document}